\documentclass[11pt,a4paper]{article}

\usepackage[margin=2.4cm]{geometry}
\usepackage[T1]{fontenc}
\usepackage[utf8]{inputenc}
\usepackage{lmodern}
\usepackage{microtype}
\usepackage{enumitem}
\usepackage{xcolor}
\usepackage[hidelinks]{hyperref}

\setlist[enumerate]{itemsep=0.45em, topsep=0.5em}
\setlist[itemize]{itemsep=0.3em, topsep=0.35em}

\newcommand{\menu}[1]{\textbf{#1}}
\newcommand{\file}[1]{\texttt{#1}}
\newcommand{\cmd}[1]{\texttt{#1}}

\title{Before the first session\\
\large From an empty computer to the course folder}
\author{AI and Data Science}
%\date{Read this before the course. Do not open the course README yet.}

\begin{document}
\maketitle

\begin{quote}
\noindent
This sheet helps install the tools used in the course. It stops when the course folder is on your computer. The instructor will then explain every cell.
\end{quote}

\tableofcontents
\vspace{0.8em}

\section{What you are installing, and when to stop}

You need five things.

\begin{itemize}
\item A \textbf{GitHub account}: the website that stores the course.
\item \textbf{Python~3.12}: the language used in the cells. Install 3.12, not a newer 3.14.
\item \textbf{Visual Studio Code}: the program in which you will open the chapters.
\item The \textbf{Jupyter extension} inside Visual Studio Code: the piece that runs one cell at a time.
\item \textbf{Git}: the tool that copies the course from GitHub onto your computer. The copy is called a \textbf{clone}.
\end{itemize}

At the end you have a folder named \file{data-extraction-lesson} on your computer, with a file \file{README.md} inside it. Stop there. Do not install the course packages. Do not run a cell. The instructor starts from the README.

Use a laptop you can bring to class, with a browser and the right to install programs. Plug it in. The downloads are a few hundred megabytes.

\section{Create the GitHub account}
\label{sec:github}

Do this first. The confirmation email can arrive while the other programs install.

\begin{enumerate}
\item Open \url{https://github.com/signup}.
\item Enter an email address you can open now. Choose a password. Choose a username. The username is public. Write it down: the instructor may need it.
\item Choose the free account. You do not need GitHub Pro.
\item Open the confirmation message and confirm the address.
\item Sign in at \url{https://github.com/login} and check that you see your home page.
\end{enumerate}

If the course repository is still private, send that username to the instructor before the session. A private page returns ``404'' until you have been added.

\section{Install Python 3.12}

Python is the language of the course. The chapters were prepared with Python~3.12. A newer Python~3.14 can be installed beside it and then be selected by mistake. Install 3.12 and use that one.

The large download button on \url{https://www.python.org/downloads/} is a newer Python. This course uses \textbf{Python 3.12.10}, the last 3.12 release that still provides a Mac installer and a Windows installer. Open this page and use only the files named there:

\begin{quote}
\url{https://www.python.org/downloads/release/python-31210/}
\end{quote}

\subsection{Mac}

\begin{enumerate}
\item On the Python~3.12 release page, under \menu{Files}, download the \textbf{macOS 64-bit universal2 installer} (\file{.pkg}).
\item Open the downloaded file. Click \menu{Continue}, then \menu{Install}. Approve with your Mac password.
\item If macOS says the file is from an unidentified developer: open \menu{System Settings}, then \menu{Privacy \& Security}, and click \menu{Open Anyway}.
\item Open \menu{Terminal}. It is in \menu{Applications}, then \menu{Utilities}.
\item Type the following line and press Return.
\end{enumerate}

\begin{verbatim}
python3 --version
\end{verbatim}

The answer must contain \cmd{3.12}. An answer that contains \cmd{3.14}, or a window that offers to install ``command line developer tools'', means another Python answered. Close that window. The python.org install is the one this course uses.

\subsection{Windows}

\begin{enumerate}
\item On the Python~3.12 release page, under \menu{Files}, download \textbf{Windows installer (64-bit)}.
\item Open the downloaded file. On the first screen, tick \menu{Add python.exe to PATH}. This box is at the bottom. The course will not find Python if it stays empty.
\item Click \menu{Install Now}. Approve the security question.
\item If the installer offers to disable the path length limit, accept it.
\item Open \menu{Command Prompt}: click the Start menu, type \cmd{cmd}, and press Enter.
\item Type the following line and press Enter.
\end{enumerate}

\begin{verbatim}
python --version
\end{verbatim}

The answer must contain \cmd{3.12}. If the window says that \cmd{python} is not recognized, the PATH box was not ticked. Run the installer again, choose \menu{Modify}, and enable \menu{Add Python to environment variables}. Then close Command Prompt, open it again, and repeat the version command.

\section{Install Visual Studio Code}

Visual Studio Code, often called VS Code, is the editor. You will open the chapter files there.

\begin{enumerate}
\item Open \url{https://code.visualstudio.com/Download}.
\item Download the installer for your computer: \menu{Mac} (Apple silicon if the Mac is from about 2020 onward, Intel if it is older), or \menu{Windows} \menu{User Installer}, 64 bit.
\item Mac: open the \file{.zip}, and drag \menu{Visual Studio Code} into \menu{Applications}. Open it from \menu{Applications}. Windows: run the installer and accept the suggested options. Tick \menu{Add to PATH} if you see it.
\item The first launch asks about theme and language. The defaults are fine. You can close the welcome tips.
\end{enumerate}

\section{Install Jupyter inside Visual Studio Code}

Jupyter, for this course, is not a second program. It is an extension of VS Code. It is what places a result under a cell when the instructor asks you to run that cell. You do not need Anaconda, and you do not need a separate JupyterLab.

\begin{enumerate}
\item In VS Code, click the extensions icon in the left bar. It looks like four squares. You can also open the menu \menu{View}, then \menu{Extensions}.
\item In the search box, type \cmd{Python}. Choose the extension named \textbf{Python} whose publisher is \textbf{Microsoft}. Click \menu{Install}.
\item Search for \cmd{Jupyter}. Choose the extension named \textbf{Jupyter} whose publisher is \textbf{Microsoft}. Click \menu{Install}.
\item Wait until both buttons say \menu{Installed} or \menu{Disable}. A \menu{Reload} button, if it appears, should be clicked.
\end{enumerate}

Leave VS Code open. The next section uses it.

\section{Install Git}

Git is the program that downloads the course and keeps the link with GitHub. The download itself is called a clone.

\subsection{Mac}

\begin{enumerate}
\item Open \menu{Terminal}.
\item Type the following line and press Return.
\end{enumerate}

\begin{verbatim}
git --version
\end{verbatim}

\begin{enumerate}
\setcounter{enumi}{2}
\item If a version number appears, Git is already there. Go to Section~\ref{sec:clone}.
\item If macOS offers to install command line developer tools, click \menu{Install} and wait. It can take several minutes. Then run \cmd{git --version} again.
\item If that offer does not appear, download the macOS installer from \url{https://git-scm.com/download/mac} and run it. Then run \cmd{git --version} again.
\end{enumerate}

\subsection{Windows}

\begin{enumerate}
\item Open \url{https://git-scm.com/download/win}. The download of \textbf{64-bit Git for Windows} starts from that page.
\item Run the installer. Keep the suggested options. One screen is titled ``Adjusting your PATH environment''. Leave the selected choice as it is: Git will be available from the command line.
\item Open a new Command Prompt and type:
\end{enumerate}

\begin{verbatim}
git --version
\end{verbatim}

A version number means Git is installed. If the command is not recognized, close Command Prompt, open it again, and retry. A window that was already open does not see a program installed a moment ago.

\section{Clone the course folder}
\label{sec:clone}

Cloning creates the folder \file{data-extraction-lesson} on your computer and links it to GitHub. Put it in \menu{Documents}, not on a shared drive and not inside another course folder.

The address to copy is exactly this one:

\begin{quote}
\url{https://github.com/cheuchenne/data-extraction-lesson.git}
\end{quote}

\begin{enumerate}
\item In VS Code, open the Source Control icon in the left bar. It looks like a branch. You can also use the menu \menu{View}, then \menu{Source Control}.
\item Click \menu{Clone Repository}. If you do not see that button, press \cmd{Ctrl+Shift+P} on Windows or \cmd{Command+Shift+P} on a Mac, type \cmd{Git: Clone}, and press Enter.
\item Paste the address above and confirm.
\item Choose \menu{Documents} as the place to save the folder. Do not open a folder that already contains files of yours.
\item When VS Code asks whether to open the cloned repository, click \menu{Open}.
\item The first time, a browser window asks you to sign in to GitHub. Sign in with the account from Section~\ref{sec:github} and authorize Visual Studio Code.
\end{enumerate}

\subsection{If VS Code does not offer Clone}

Use the terminal you already opened. On a Mac, in Terminal:

\begin{verbatim}
cd ~/Documents
git clone https://github.com/cheuchenne/data-extraction-lesson.git
\end{verbatim}

On Windows, in Command Prompt:

\begin{verbatim}
cd %USERPROFILE%\Documents
git clone https://github.com/cheuchenne/data-extraction-lesson.git
\end{verbatim}

Then, in VS Code, use \menu{File}, then \menu{Open Folder}, and choose the new folder \file{data-extraction-lesson} inside \menu{Documents}.

\section{Check the folder, then stop}

In the left-hand file list of VS Code you should see, among other names:

\begin{itemize}
\item \file{README.md}
\item a folder \file{lessons}
\item a folder \file{python}
\item a folder \file{data}
\item a file \file{requirements.txt}
\end{itemize}

That is the end of this sheet. Do not follow the installation commands inside \file{README.md} yet. Do not run a cell. Those steps are the start of the session, and the instructor will do them with you.

\section{If a step fails}

\begin{description}
\item[The GitHub page says 404.] The repository is private and your username has not been added. Send the username from Section~\ref{sec:github} to the instructor. Do not create a new repository of your own.
\item[The clone asks for a password and then refuses it.] GitHub does not accept the account password in that box. In VS Code, use the browser sign-in from Section~\ref{sec:clone}. If the terminal is asking, stop and use \menu{Git: Clone} in VS Code instead.
\item[\cmd{python} or \cmd{python3} is not 3.12.] The computer is answering with another installation. Tell the instructor the exact line the version command printed. Do not uninstall anything before the session.
\item[VS Code has no cell buttons.] The Jupyter extension from the Microsoft publisher is missing or not yet enabled. Return to the extensions screen and install it.
\item[The download is blocked by the school network.] Try another network, or download at home. The same installers are the ones linked above.
\end{description}

\end{document}
